\documentclass[10pt,twoside]{article}\usepackage[]{graphicx}\usepackage[dvipsnames,svgnames,table]{xcolor}
% maxwidth is the original width if it is less than linewidth
% otherwise use linewidth (to make sure the graphics do not exceed the margin)
\makeatletter
\def\maxwidth{ %
  \ifdim\Gin@nat@width>\linewidth
    \linewidth
  \else
    \Gin@nat@width
  \fi
}
\makeatother

\definecolor{fgcolor}{rgb}{0.345, 0.345, 0.345}
\newcommand{\hlnum}[1]{\textcolor[rgb]{0.686,0.059,0.569}{#1}}%
\newcommand{\hlsng}[1]{\textcolor[rgb]{0.192,0.494,0.8}{#1}}%
\newcommand{\hlcom}[1]{\textcolor[rgb]{0.678,0.584,0.686}{\textit{#1}}}%
\newcommand{\hlopt}[1]{\textcolor[rgb]{0,0,0}{#1}}%
\newcommand{\hldef}[1]{\textcolor[rgb]{0.345,0.345,0.345}{#1}}%
\newcommand{\hlkwa}[1]{\textcolor[rgb]{0.161,0.373,0.58}{\textbf{#1}}}%
\newcommand{\hlkwb}[1]{\textcolor[rgb]{0.69,0.353,0.396}{#1}}%
\newcommand{\hlkwc}[1]{\textcolor[rgb]{0.333,0.667,0.333}{#1}}%
\newcommand{\hlkwd}[1]{\textcolor[rgb]{0.737,0.353,0.396}{\textbf{#1}}}%
\let\hlipl\hlkwb

\usepackage{framed}
\makeatletter
\newenvironment{kframe}{%
 \def\at@end@of@kframe{}%
 \ifinner\ifhmode%
  \def\at@end@of@kframe{\end{minipage}}%
  \begin{minipage}{\columnwidth}%
 \fi\fi%
 \def\FrameCommand##1{\hskip\@totalleftmargin \hskip-\fboxsep
 \colorbox{shadecolor}{##1}\hskip-\fboxsep
     % There is no \\@totalrightmargin, so:
     \hskip-\linewidth \hskip-\@totalleftmargin \hskip\columnwidth}%
 \MakeFramed {\advance\hsize-\width
   \@totalleftmargin\z@ \linewidth\hsize
   \@setminipage}}%
 {\par\unskip\endMakeFramed%
 \at@end@of@kframe}
\makeatother

\definecolor{shadecolor}{rgb}{.97, .97, .97}
\definecolor{messagecolor}{rgb}{0, 0, 0}
\definecolor{warningcolor}{rgb}{1, 0, 1}
\definecolor{errorcolor}{rgb}{1, 0, 0}
\newenvironment{knitrout}{}{} % an empty environment to be redefined in TeX

\usepackage{alltt}
\usepackage[marginparsep=1em]{geometry}
\geometry{lmargin=1.0in,rmargin=1.0in, bmargin=1.2in,  tmargin=1.2in}
\usepackage[dvipsnames,svgnames,table]{xcolor}
\usepackage{graphicx}
\usepackage{amssymb}
\usepackage{epstopdf}
\usepackage{verbatim}
\usepackage{enumerate}
\usepackage{bm}
\usepackage{amsthm}
\usepackage{float}
\usepackage{amsmath}
\usepackage{fancyhdr}
\usepackage{hyperref}
\usepackage{mathtools}

%%%%  SHORTCUT COMMANDS  %%%%
\newcommand{\ds}{\displaystyle}
\newcommand{\Z}{\mathbb{Z}}
\newcommand{\T}{\mathcal{T}}
\newcommand{\arc}{\rightarrow}
\newcommand{\R}{\mathbb{R}}
\newcommand{\RP}{\mathbb{R}(+)}
\newcommand{\Rs}{\mathbb{R}^{**}}
\newcommand{\C}{\mathbb{C}}
\newcommand{\E}{\mathbb{E}}
\newcommand{\B}{\mathcal{B}}
\newcommand{\PX}{\mathcal{P}(X)}
\newcommand*\rot{\rotatebox{90}}
\newcommand*\OK{\ding{51}}
\newcommand{\N}{\mathbb{N}}
\newcommand{\Q}{\mathbb{Q}}
\newcommand{\Answer}{\vspace{2mm}\textbf{\underline{Answer}}\\}
\newcolumntype{L}{>{$}l<{$}}
\newcolumntype{C}{>{$}c<{$}}
\newcommand{\GL}{\text{GL}}
\newcommand{\SL}{\text{SL}}
\newcommand{\Var}{\text{Var}}

\newcommand{\stirling}[2]{\genfrac{\{}{\}}{0pt}{}{#1}{#2}}

\pagestyle{fancy}
\fancyhf{}
\renewcommand{\sectionmark}[1]{\markright{\thesection.\ #1}}
\lhead{\fancyplain{}{}}
\fancyhead[RE,RO]{Math 4451, Spring 2026}
\fancyfoot[RE,RO]{\thepage}
\IfFileExists{upquote.sty}{\usepackage{upquote}}{}
\begin{document}
\begin{flushright}
\begin{minipage}{.33\textwidth}
\rightline{YOUR NAME}
\rightline{\href{mailto:YOUR EMAIL}{YOUR EMAIL}}
\rightline{\today}
\end{minipage}
\end{flushright}

\begin{center}
{\large{\textbf{Homework 8}}}
\end{center}

\begin{enumerate}
  \item While UMVUE estimators are desirable, there are not always available.
  Instead, we can sometimes further restrict our class of estimators to unbiased, and linear.
  In this class of estimators, we often what the Best Linear Unbiased Estimator (BLUE) of a parameter $\theta$.
  That is, among all linear, unbiased estimators, we want the estimator with minimum variance.

  Suppose you are measure the concentration $\theta$ of a chemical in a water sample.
  You have two machines to measure the chemical, and both machines are \textbf{unbiased}.
  If $X_1, X_2$ represent measurements from machine 1 and machine 2, respectively, then $E[X_1] = \theta$, and $E[X_2] = \theta$.
  However, the machines have different variance in their estimates: $\Var(X_1) = \sigma^2_1$, and $\Var(X_2) = \sigma^2_2$.

  There are three different proposals of how to estimate $\theta$, the true unknown chemical concentration:
  \begin{enumerate}[(A)]
    \item $\hat{\theta}_A = 0.5 X_1 + 0.5 X_2$.
    \item $\hat{\theta}_B = 0.2 X_1 + 0.8 X_2$.
    \item $\hat{\theta}_C = 0.5 X_1 + 0.4 X_2$.
  \end{enumerate}

\begin{enumerate}
  \item Which of the three estimators are \textbf{Unbiased}?
  \item If $\sigma^2_1 = \sigma^2_2$, which of the three estimators has the lowest variance? What about lowest variance amongst only the unbiased estimators?
  \item Consider any \emph{linear} estimator, $\hat{\mu} = aX_1 + bX_2$. Find a rule for $a$ and $b$ such that $\hat{\mu}$ is unbiased.
  \item Find the BLUE of $\theta$ using $X_1$ and $X_2$. That is, the linear estimator of $\theta$ that is unbiased and minimizes the variance of the estimator (hint: since we want small variance, restrict yourselves to values $0 \leq a, b \leq 1$).
\end{enumerate}

\item A factory produces computer chips. Each chip has a probability $\theta$ of being defective. The factory manager wants to know the probability of drawing \textbf{two defective chips in a row}.
Assuming the chip quality is independent across manufactured chips, the parameter we want to estimate is $\theta = p^2$.

A natural model for this scenario is that $X_1, X_2, \ldots, X_n$ are iid Bernoulli$(p)$.

\begin{enumerate}
  \item \label{prob:binomMLE} Use the invariance property of the MLE to find the MLE of $\theta = p^2$.
  \item Calculate the bias of the MLE estimate found in Problem~\ref{prob:binomMLE}.
  \item Consider the alternative estimator $\hat{\theta}_2 = X_1X_2$ (the product of the first two observations). Show that $\hat{\theta}_2$ is unbiased for $\theta = p^2$.
  \item Because the $\hat{\theta}_2$ only depends on the first two observations, it's not a very good estimator. We want to improve this using the \textbf{Rao-Blackwell Theorem}. Find a sufficient statistic, $T$ for $p$ (note that if $T$ is sufficient for $p$, it is sufficient for $p^2$). (HINT: there are many possible solutions. You'll want to simplify as much as possible to solve Problem~\ref{prob:RB}).
  \item \label{prob:RB} With your sufficient statistic $T$, apply the Rao-Blackwell theorem to find an improved unbiased estimator:
  $$
  \hat{\theta}_{RB} = E\big[ \hat{\theta}_2 | T\big] = E\big[ X_1X_2 | T\big].
  $$
  This will involve some counting practice from last semester. In particular, $X_1X_2$ is either $0$ or $1$. Therefore, $E\big[ X_1X_2 | T\big] = P(X_1X_2 = 1 | T) = P(X_1 = 1, X_2 = 1 | T)$. Thus, if we know $T = t$, what is the probability that the first two observations are equal to $1$?
\end{enumerate}


\end{enumerate}


\end{document}
